\chapter[\texorpdfstring{Algèbre\hspace{0.28em}linéaire}{Algèbre linéaire}]{Algèbre linéaire}\label{ch:iv1:linear-algebra}

\begin{ivtrains}
Lire le rang et les valeurs propres dans la structure d'une matrice, sans développement à la main, et savoir ce qu'une matrice de covariance peut et ne peut pas être. Les questions suivantes portent sur l'ACP et la SVD, sur le conditionnement, et sur la raison pour laquelle une matrice de covariance empirique induit en erreur lorsque les actifs sont nombreux et les jours peu nombreux.
\end{ivtrains}

\begin{question}[$\star$ \iqroles{QR, QT, DS, QD} \iqfirm{toutes} \iqlevel{junior}]\label{iq:iv1:linear-algebra:1}
Soient $u = (1, 2, 3)$ et $v = (1, 1, 1)$. Quels sont le rang et les valeurs propres de la matrice $uv^\top$~? Quelles sont les valeurs propres de la matrice $n$ par $n$ dont tous les coefficients valent un, et de celle dont la diagonale est nulle et dont les autres coefficients valent un~?
\end{question}

\begin{question}[$\star$ \iqroles{QR, QT, DS, Risque} \iqfirm{toutes} \iqlevel{junior}]\label{iq:iv1:linear-algebra:2}
Dix actifs ont une variance unitaire et une corrélation deux à deux commune de $\rho = 0.2$. Quelles sont les valeurs propres de la matrice de corrélation, et quels sont ses vecteurs propres~? Quelle est la plus petite valeur possible de $\rho$ pour 10 actifs, et quelle est la variance de leur moyenne à pondérations égales lorsque le nombre d'actifs est grand~?
\end{question}

\begin{question}[$\star$ \iqroles{QR, QT, Risque, DS} \iqfirm{toutes} \iqlevel{junior}]\label{iq:iv1:linear-algebra:3}
Trois actifs ont des corrélations deux à deux $\rho_{12} = 0.9$, $\rho_{23} = 0.9$ et $\rho_{13} = 0$, telles que les fournit un fournisseur de données. Est-ce une matrice de corrélation valide~? Démontrez votre réponse en trouvant un portefeuille à la variance impossible.
\end{question}

\begin{question}[$\star$ \iqroles{QR, QT, DS, QD} \iqfirm{toutes} \iqlevel{junior}]\label{iq:iv1:linear-algebra:4}
Une matrice $P$ vérifie $P^2 = P$. Quelles peuvent être ses valeurs propres~? Montrez que son rang est égal à sa trace. Quelle est la trace de la matrice $X(X^\top X)^{-1}X^\top$, où $X$ a $n$ lignes et $p$ colonnes indépendantes~?
\end{question}

\begin{question}[$\star$ \iqroles{QR, QT, QD, DS} \iqfirm{toutes} \iqlevel{junior}]\label{iq:iv1:linear-algebra:5}
Soit $A$ la matrice symétrique 2 par 2 dont les coefficients diagonaux valent 2 et dont les autres coefficients valent 1. Trouvez ses valeurs propres et ses vecteurs propres, et utilisez-les pour écrire exactement $A^{10}$ et $A^{-1}$, sans multiplication répétée.
\end{question}

\begin{question}[$\star\star$ \iqroles{QR, QT, Risque, DS} \iqfirm{toutes} \iqlevel{junior}]\label{iq:iv1:linear-algebra:6}
Un actif $X$ présente une corrélation de $0.9$ avec $Y$, et $Y$ présente une corrélation de $0.9$ avec $Z$. Quelles sont les corrélations minimale et maximale possibles entre $X$ et $Z$~? Déterminez-les de deux façons~: à partir d'un déterminant, puis à partir des angles entre vecteurs.
\end{question}

\begin{question}[$\star\star$ \iqroles{QR, QT, QD} \iqfirm{toutes} \iqlevel{junior}]\label{iq:iv1:linear-algebra:7}
Soit $P$ le décalage cyclique $n$ par $n$, avec $(Px)_i = x_{i + 1}$, les indices se refermant sur eux-mêmes. Quelles sont ses valeurs propres et ses vecteurs propres, et quel est son conditionnement~? Utilisez-les pour trouver les valeurs propres de la matrice circulante 4 par 4 dont la première ligne est $(2, 1, 0, 0)$.
\end{question}

\begin{question}[$\star\star$ \iqroles{QR, QD, QT} \iqfirm{toutes} \iqlevel{junior}]\label{iq:iv1:linear-algebra:8}
Prenez la matrice tridiagonale $n$ par $n$ dont la diagonale vaut 2 et dont les coefficients voisins de la diagonale valent $-1$. Trouvez ses valeurs propres pour $n = 3$, puis donnez la formule générale. Que devient le conditionnement lorsque $n$ augmente, et pourquoi cela importe-t-il pour résoudre un problème aux différences finies~?
\end{question}

\begin{question}[$\star\star$ \iqroles{QR, QT, DS, Risque} \iqfirm{hedge fund} \iqlevel{junior}]\label{iq:iv1:linear-algebra:9}
Deux actifs ont une corrélation de $0.5$. Si tous deux ont une volatilité unitaire, quelles sont les deux composantes principales et leurs parts de variance~? Si la volatilité du premier actif est doublée, quelle part de variance la première composante explique-t-elle~? Pourquoi cela montre-t-il que l'ACP dépend des unités, et que faites-vous pour y remédier~?
\end{question}

\begin{question}[$\star\star$ \iqroles{QR, QD, DS} \iqfirm{toutes} \iqlevel{junior}]\label{iq:iv1:linear-algebra:10}
Déterminez les valeurs singulières de la matrice dont les lignes sont $(3, 0)$ et $(4, 5)$, ainsi que son conditionnement pour la norme 2. Comparez-les aux valeurs propres de cette matrice. Comment les valeurs singulières se rapportent-elles aux composantes principales d'une matrice de données~?
\end{question}

\begin{question}[$\star\star$ \iqroles{QR, QD, QT, DS} \iqfirm{toutes} \iqlevel{junior}]\label{iq:iv1:linear-algebra:11}
Résolvez $Ax = b$ pour $A$ dont les lignes sont $(1, 1)$ et $(1, 1.0001)$, et pour $b = (2, 2)$ puis $b = (2, 2.0001)$. Que montre le résultat~? Donnez le conditionnement de $A$, ainsi que la règle empirique qui le relie aux chiffres perdus. Que produit l'élévation au carré de $A$, comme dans les équations normales d'une régression~?
\end{question}

\begin{question}[$\star\star$ \iqroles{QR, QT, DS, Risque} \iqfirm{hedge fund, banque} \iqlevel{junior}]\label{iq:iv1:linear-algebra:12}
Vous calculez la matrice de covariance empirique de 500 actions à partir de 250 rendements quotidiens. Quel est le plus grand rang qu'elle peut atteindre~? Qu'en découle-t-il pour un portefeuille de variance minimale calculé à partir d'elle~? Donnez deux remèdes pratiques.
\end{question}

\begin{question}[$\star\star$ \iqroles{QR, QT, Risque, DS} \iqfirm{hedge fund} \iqlevel{junior}]\label{iq:iv1:linear-algebra:13}
Une matrice de covariance a pour valeurs propres 5, 1 et $0.1$. Vous la remplacez par $0.7$ fois elle-même plus $0.3$ fois $\mu I$, où $\mu$ est la moyenne des valeurs propres. Quelles sont les nouvelles valeurs propres et le nouveau conditionnement~? Qu'est-ce qui reste identique~? Pourquoi cela aide-t-il un optimiseur, et d'où viendrait $0.3$~?
\end{question}

\begin{question}[$\star\star$ \iqroles{QR, QT, Risque} \iqfirm{hedge fund} \iqlevel{junior}]\label{iq:iv1:linear-algebra:14}
Un actif a une volatilité de 10~\%, un autre de 20~\%, et leur corrélation vaut $\rho$. Déterminez les pondérations de variance minimale. Pour quelle valeur de $\rho$ le poids du second actif est-il exactement nul~? Quels sont les poids et la volatilité du portefeuille pour $\rho = 0.8$~? Pourquoi la variance est-elle inférieure à celle du premier actif, alors que le second est plus risqué et corrélé~?
\end{question}

\begin{question}[$\star\star$ \iqroles{QR, QT, DS, QD} \iqfirm{toutes} \iqlevel{junior}]\label{iq:iv1:linear-algebra:15}
Un état de marché passe d'un régime à l'autre~: depuis le régime calme, il reste calme avec une probabilité de $0.9$~; depuis le régime tendu, il reste tendu avec une probabilité de $0.8$. Écrivez la matrice de transition, ses valeurs propres et sa loi stationnaire. Avec quelle rapidité l'effet de l'état d'aujourd'hui s'estompe-t-il~? Quelle est cette durée pour une décroissance de moitié~?
\end{question}

\begin{question}[$\star\star$ \iqroles{QR, QD, Dev} \iqfirm{toutes} \iqlevel{junior}]\label{iq:iv1:linear-algebra:16}
Écrivez la récurrence de Fibonacci sous forme d'une matrice 2 par 2. Quelles sont ses valeurs propres~? Dérivez la forme close, et dites combien de multiplications les puissances matricielles nécessitent pour $F_n$. En utilisant l'arithmétique à virgule flottante en double précision, la forme close donne-t-elle toujours l'entier exact~?
\end{question}

\begin{question}[$\star\star$ \iqroles{QR, QT, QD} \iqfirm{toutes} \iqlevel{junior}]\label{iq:iv1:linear-algebra:17}
Une matrice 100 par 100 a pour coefficients $a_{ij} = i + j$, les indices allant de 1 à 100. Quel est son rang~? Trouvez ses valeurs propres non nulles sous forme close, sans rien développer, en utilisant la structure de la matrice.
\end{question}

\begin{question}[$\star\star\star$ \iqroles{QR, QT, DS, Risque} \iqfirm{hedge fund} \iqlevel{junior, confirmé}]\label{iq:iv1:linear-algebra:18}
Cinq cents actifs non corrélés, de variance unitaire, sont observés pendant 1000 jours. Entre quelles valeurs se situent les valeurs propres de la matrice de covariance empirique, selon la loi de Marchenko--Pastur, alors que toutes les vraies valeurs propres valent 1~? Un portefeuille de variance minimale est construit à partir de la matrice de covariance empirique. Comment son risque dans l'échantillon (in-sample) se compare-t-il à son risque hors échantillon (out-of-sample), et à l'optimum véritable~?
\end{question}

\begin{question}[$\star\star\star$ \iqroles{QR, QT, QD, Risque} \iqfirm{hedge fund} \iqlevel{junior, confirmé}]\label{iq:iv1:linear-algebra:19}
Un modèle factoriel de risque a pour matrice de covariance $\Sigma = D + \beta\beta^\top$ pour 1000 actifs, avec $D$ diagonale. Déterminez $\Sigma^{-1}$ avec la formule de Sherman--Morrison, et démontrez-la. Quel est le coût par rapport à une inversion directe~? Écrivez les pondérations de variance minimale, et dites ce qu'elles font aux actions à bêta élevé.
\end{question}

\begin{question}[$\star\star\star$ \iqroles{QR, QT, DS} \iqfirm{hedge fund} \iqlevel{junior, confirmé}]\label{iq:iv1:linear-algebra:20}
$X$, $Y$ et $Z$ ont une variance unitaire, et toutes leurs corrélations deux à deux valent $0.5$. Quelle est la corrélation partielle de $X$ et $Y$ sachant $Z$~? Calculez-la de deux façons~: à partir des résidus des régressions sur $Z$, puis à partir de l'inverse de la matrice de corrélation. Quelle est la corrélation partielle en général si la corrélation vaut $\rho$~?
\end{question}

\begin{question}[$\star\star\star$ \iqroles{QR, QD, DS} \iqfirm{toutes} \iqlevel{junior, confirmé}]\label{iq:iv1:linear-algebra:21}
Vous avez besoin du vecteur propre dominant d'une matrice de covariance 1000 par 1000, dont les deux plus grandes valeurs propres valent 10 et 5. Comment fonctionne la méthode de la puissance itérée, et combien d'itérations faut-il environ pour obtenir une erreur de $10^{-6}$ sur le vecteur~? Et sur la valeur propre~? Combien d'itérations faut-il si la deuxième valeur propre vaut $9.9$~? Dans quels cas pourrait-elle échouer~?
\end{question}

\begin{question}[$\star\star\star$ \iqroles{QR, QD, DS, QT} \iqfirm{hedge fund} \iqlevel{junior, confirmé}]\label{iq:iv1:linear-algebra:22}
Montrez que $AB$ et $BA$ ont les mêmes valeurs propres non nulles, même si leurs tailles diffèrent. Utilisez ce résultat pour calculer les composantes principales d'une matrice de rendements $T$ par $N$ avec $T = 250$ et $N = 5000$, sans former la matrice $N$ par $N$. Quelle est l'application qui relie les vecteurs propres~? Dérivez aussi $\det(I + uv^\top) = 1 + v^\top u$.
\end{question}

\begin{question}[$\star\star\star$ \iqroles{QR, QT, Risque, QD} \iqfirm{banque, hedge fund} \iqlevel{junior, confirmé}]\label{iq:iv1:linear-algebra:23}
Une matrice de corrélation construite à partir d'un mélange de sources a $\rho_{12} = \rho_{23} = 0.9$ et $\rho_{13} = 0$. Donnez ses valeurs propres. Réparez-la en écrêtant la valeur propre négative à zéro, puis en la remettant à l'échelle pour obtenir une diagonale unitaire, et donnez les nouvelles corrélations. En quel sens est-elle proche de la matrice d'origine, et est-ce la plus proche possible~?
\end{question}

\begin{question}[$\star\star\star\star$ \iqroles{QR, QT} \iqfirm{hedge fund} \iqlevel{confirmé}]\label{iq:iv1:linear-algebra:24}
Une matrice $n$ par $n$ a des coefficients indépendants, chacun d'espérance nulle et de variance 1. Déterminez $E[\det(A)^2]$ exactement. Vérifiez la réponse pour $n = 3$ avec des coefficients $\pm 1$, en raisonnant sur la taille de l'énumération. Qu'est-ce que cela vous apprend sur la façon dont la taille typique d'un déterminant croît~? Comparez-la au plus grand déterminant possible d'une matrice à coefficients $\pm 1$.
\end{question}

\begin{question}[$\star\star\star\star$ \iqroles{QR, QT, QD} \iqfirm{hedge fund} \iqlevel{confirmé}]\label{iq:iv1:linear-algebra:25}
Démontrez que le produit terme à terme de deux matrices semi-définies positives est lui-même semi-défini positif. Utilisez ce résultat pour montrer que le produit d'une matrice de corrélation, terme à terme, par une matrice de coefficients $e^{-\lambda |i - j|}$ donne une matrice de corrélation valide. Démontrez aussi la borne $\lambda_{\min}(A \circ B) \ge \lambda_{\min}(A) \min_i B_{ii}$.
\end{question}
